\documentclass[11pt]{article}
\usepackage[margin=1in]{geometry}
\usepackage[T1]{fontenc}
\usepackage[utf8]{inputenc}
\usepackage{lmodern}
\usepackage{amsmath,amssymb,mathtools,booktabs,graphicx,natbib,hyperref}
\usepackage[nameinlink,noabbrev]{cleveref}

\title{Numerical Selection of Forecast-Optimal Smoothness in Penalized Trend Estimation}
\author{Heriberto Espino Montelongo}
\date{\today}
\newcommand{\R}{\mathbb{R}}

\begin{document}
\maketitle

\begin{abstract}
We study numerical selection of normalized smoothness for finite-difference
penalized trend estimation when the objective is chronological forecast loss.
The penalty parameter is reparameterized by a normalized smoothness index on
the compact domain $[0,1]$. Because the forecast-validation objective need not
be unimodal, the numerical task is to discover, refine, and classify multiple
stationary points without relying on an exhaustive dense grid. The proposed
method combines a sparse deterministic partition, derivative-aware adaptive
subdivision, bracketed root refinement, and explicit comparison of local
minima and boundaries. A dense reference search is used only to quantify
optimum agreement, missed-minimum risk, objective regret, evaluation count,
and runtime. This manuscript remains a research draft until the numerical
benchmark protocol is frozen.
\end{abstract}

\section{Problem formulation}

For $y\in\R^N$, let $D_d$ be the order-$d$ finite-difference matrix and
$Q=D_d^\top D_d$. The penalized trend is
\[
\widehat\tau_\lambda
=
\arg\min_{\tau\in\R^N}
\left\{
\|y-\tau\|_2^2+\lambda\|D_d\tau\|_2^2
\right\}
=
(I+\lambda Q)^{-1}y.
\]

\section{Normalized smoothness}

Let $\delta_1,\ldots,\delta_{N-d}$ be the positive eigenvalues of $Q$.
Define
\[
S(\lambda)
=
1-
\frac{1}{N-d}
\sum_{j=1}^{N-d}
\frac{1}{1+\lambda\delta_j}.
\]
Then $S(0)=0$, $S(\lambda)\to1$ as $\lambda\to\infty$, and
\[
S'(\lambda)
=
\frac{1}{N-d}
\sum_{j=1}^{N-d}
\frac{\delta_j}{(1+\lambda\delta_j)^2}>0.
\]

\section{Forecast-validation objective}

For fixed discrete choices $(d,L,m,h)$, where $m$ is a pre-specified trend
continuation rule, let
\[
F(S)=CV_h(d,L,m,S).
\]
The continuous optimization problem is
\[
S^\star\in\arg\min_{S\in[0,1]}F(S).
\]

If $f(\lambda)=F(S(\lambda))$, then
\[
F'(S)=\frac{f'(\lambda)}{S'(\lambda)}.
\]
Hence interior stationary points are preserved by the monotone
reparameterization.

\section{Multiple-stationary-point search}

The proposed method begins from a small deterministic partition of the
smoothness domain, evaluates derivative information, adaptively subdivides
intervals with unresolved stationary structure, brackets derivative sign
changes, refines roots with a bracketed solver, classifies stationary points,
and compares the detected local minima with the exact or limiting boundaries.

Nearby local minima may be summarized after discovery with an
$\varepsilon$-spacing rule. This post-processing step does not alter the root
search itself.

\section{Numerical validation}

The adaptive method is compared with a very dense smoothness reference. Primary
metrics are selected-smoothness error, objective regret, local minima
found/missed, evaluation count, runtime, and failure or ambiguity rate.
Controlled experiments vary difference order, estimation-window length,
forecast horizon, and a small set of simple continuation rules only to create
different objective geometries.

\section{Scope}

The paper does not study adaptive time-varying selection of $(d,L,S)$ and does
not perform a broad financial comparison among ARIMA, state-space, likelihood,
GCV, or recurrence models. Those are separate research tracks in the same
repository.

\section{Limitations}

Without stronger assumptions, finite adaptive sampling cannot certify discovery
of every stationary point of an arbitrary smooth objective. The empirical
benchmark therefore measures missed-root and missed-minimum behavior rather
than claiming a universal guarantee.

\bibliographystyle{plainnat}
\IfFileExists{literature/references.bib}{
  \bibliography{literature/references}
}{
  \bibliography{../literature/references}
}
\end{document}
